\section{Proposed study design}
\label{sec:plan}

The literature review points to one gap more than any
other: food-intake detectors from wearable chewing sensors are usually
reported as a single architecture with a single number, and the design
decisions behind them are rarely isolated. The ablation in
Table~\ref{tab:ablation} already shows that those decisions are not equally
important, since the input representation accounts for most of the
performance and the attention components cannot be separated from noise.
The proposed study makes that the question rather than a side result: which
design dimensions matter for deep food-intake detection from a chewing
sensor, measured one at a time under leave-one-subject-out validation.

\begin{table}[H]
\centering
\caption{Preliminary ablation: the model with one component removed or
added. The first 20 annotated sessions, participant-level 4-fold
cross-validation, same folds for every row. Fold-to-fold spread in F1 was
about 0.048, so effects smaller than that are not distinguishable from
noise. To be repeated on all 61 sessions.}
\label{tab:ablation}
\begin{tabular}{@{}lrrrl@{}}
\toprule
\rowcolor{black!8}
\textbf{Variant} & \textbf{F1} & \textbf{Change} & \textbf{Precision} & \textbf{Verdict} \\
\midrule
Full model                             & 0.854 & --     & 0.876 & -- \\
FFT magnitude input instead of time series & 0.767 & −0.087 & 0.696 & Hurts \\
No overlapping training windows        & 0.829 & −0.025 & 0.834 & Possibly hurts \\
No attention (global average pooling)  & 0.847 & −0.007 & 0.857 & Within noise \\
No chew-count head                     & 0.849 & −0.006 & 0.857 & Within noise \\
Add median smoothing (3 windows)       & 0.775 & −0.079 & 0.765 & Hurts \\
Add 8\,s context each side, no marker  & 0.775 & −0.080 & 0.683 & Hurts \\
\bottomrule
\end{tabular}
\end{table}

\subsection{Protocol, fixed for every experiment}

\begin{itemize}
  \item \textbf{Leave-one-subject-out validation} on all 43 participants
        (43 folds, each holding out every session of one participant),
        following \citet{saeb2017need}. At about a minute per fold this is
        computationally feasible, and it is the closest approximation of
        the use case: a participant the model has never seen. Five-fold
        runs are used only to explore.
  \item \textbf{Paired comparison per participant.} Each variant is scored
        on the same held-out participants as the reference model, so the
        difference can be tested per participant (Wilcoxon signed-rank)
        rather than read against the spread of fold means.
  \item \textbf{Three seeds per configuration}, reported as mean and range,
        so initialisation noise is not mistaken for an effect.
  \item \textbf{False positives per hour of non-eating} reported alongside
        F1, because that is the failure mode that matters in free-living use
        and the one negative data (below) is meant to fix.
\end{itemize}

\subsection{Dimensions to explore}

\begin{xltabular}{\textwidth}{@{}>{\raggedright\arraybackslash}p{27mm}
                               >{\raggedright\arraybackslash}X
                               >{\raggedright\arraybackslash}X
                               >{\raggedright\arraybackslash}p{26mm}@{}}
\caption{Planned experiments. Each changes one dimension of the model in
the model and holds everything else fixed.}
\label{tab:plan}\\
\toprule
\rowcolor{black!8}
\textbf{Dimension} & \textbf{Levels} & \textbf{Question} & \textbf{Grounding} \\
\midrule
\endfirsthead
\toprule
\rowcolor{black!8}
\textbf{Dimension} & \textbf{Levels} & \textbf{Question} & \textbf{Grounding} \\
\midrule
\endhead
\bottomrule
\endlastfoot
Published architectures & FCN, ResNet and Inception as compared by
  \citet{ghosh2022comparative}, trained on the same folds &
  Where does this model stand against the closest prior work on AIM-2, under
  an identical protocol? &
  \citet{ghosh2022comparative, wang2017strong} \\
Input representation & Time series, FFT magnitude, wavelet scalogram &
  Does keeping temporal structure explain the gain, and does a
  time-frequency image recover it? &
  \citet{ghosh2024integrated, hannun2019cardiologist} \\
Sampling rate & 128, 64, 32, 16\,Hz &
  How much of the signal is needed? Chewing lies below about 2.2\,Hz, so
  32\,Hz still leaves a Nyquist margin of seven times. &
  \citet{po2011chewing} \\
Window length & 2, 4, 8, 16\,s &
  Resolution against context: shorter windows localise bites, longer
  windows contain more chew cycles. &
  \citet{banos2014window} \\
Context & None; 8\,s each side with a target-window marker channel &
  Can surrounding signal help without the model answering "is there eating
  anywhere nearby?" &
  (this work) \\
Chew-head weight & 0, 0.1, 0.2, 0.5, 1.0 &
  How strongly should counting chews shape the features, and does it help
  participants with weak optical signal? &
  \citet{caruana1997multitask} \\
Temporal model & Attention pooling; small Transformer encoder over the 64
  convolutional time steps &
  Does self-attention between time steps add to a convolutional front end
  at this data size? &
  \citet{ikram2025transformer, zerveas2021transformer, zeng2023transformers} \\
Self-supervised pre-training & None; convolutional trunk pre-trained on
  unannotated AIM-2 recordings in ETS &
  Can the unannotated sensor data, most of what the database holds, stand
  in for labels the study does not have? &
  \citet{diou2022intake} \\
Negative data & Annotated sessions only; plus non-eating recordings from
  other ETS studies &
  Do additional negative participants reduce false positives without
  costing recall? &
  \citet{bell2020scoping, ghosh2024integrated} \\
\end{xltabular}

The published-architecture row comes first: it is the comparison a reviewer
will ask for, and it replaces comparisons with the lab's earlier model. The
representation, sampling-rate and window rows are cheap and answer the
questions raised at the 9/23 lab meeting directly. The Transformer row is the one most likely to be
asked about by reviewers, and the one most at risk of overfitting on a
dataset this size; with 24 hours of annotated data it is now worth running, and
should be reported whether it helps or not. Negative
data is expected to have the largest practical effect, since false
positives are the model's main weakness, and is the natural joint
experiment with the lab.
